\documentclass[12pt]{article}

% ============================================================
% IDIOMA Y CODIFICACIÓN
% ============================================================
\usepackage[spanish]{babel}
\usepackage[utf8]{inputenc}
\usepackage[T1]{fontenc}

% ============================================================
% MÁRGENES Y GEOMETRÍA
% ============================================================
\usepackage[a4paper, margin=2.5cm]{geometry}

% ============================================================
% ESPACIADO MEJORADO
% ============================================================
\usepackage{setspace}
\setstretch{1.15}

\setlength{\parskip}{0.5em}
\setlength{\parindent}{0pt}

% ============================================================
% TIPOGRAFÍA
% ============================================================
\usepackage{lmodern}

% ============================================================
% COMILLAS TIPOGRÁFICAS (compatibles con babel/español)
% ============================================================
\usepackage{csquotes}

% ============================================================
% FORMATO DE TÍTULOS
% ============================================================
\usepackage{titlesec}

\titleformat{\section}
  {\Large\bfseries}
  {\thesection.}
  {0.5em}
  {}

\titlespacing*{\section}{0pt}{1.5em}{0.8em}

% ============================================================
% MATEMÁTICAS
% ============================================================
\usepackage{amsmath, amssymb, amsfonts}

% ============================================================
% GRÁFICOS
% ============================================================
\usepackage{graphicx}
\usepackage{caption}

% ============================================================
% TABLAS
% ============================================================
\usepackage{booktabs}
\usepackage{array}
\usepackage{multirow}

% ============================================================
% CÓDIGO / PYTHON
% ============================================================
\usepackage{listings}
\usepackage{xcolor}

\definecolor{codebg}{RGB}{248, 248, 248}
\definecolor{codeframe}{RGB}{200, 200, 200}
\definecolor{codecomment}{RGB}{106, 153, 85}
\definecolor{codekeyword}{RGB}{86, 156, 214}
\definecolor{codestring}{RGB}{206, 145, 120}
\definecolor{codenumber}{RGB}{181, 206, 168}

\lstdefinestyle{estilo_base}{
    backgroundcolor=\color{codebg},
    frame=single,
    rulecolor=\color{codeframe},
    basicstyle=\ttfamily\small,
    commentstyle=\color{codecomment}\itshape,
    keywordstyle=\color{codekeyword}\bfseries,
    stringstyle=\color{codestring},
    numberstyle=\tiny\color{codenumber},
    numbers=left,
    numbersep=8pt,
    breaklines=true,
    breakatwhitespace=true,
    tabsize=4,
    showstringspaces=false,
    captionpos=b,
    literate=
        {á}{{\'a}}1 {é}{{\'e}}1 {í}{{\'i}}1 {ó}{{\'o}}1 {ú}{{\'u}}1
        {Á}{{\'A}}1 {É}{{\'E}}1 {Í}{{\'I}}1 {Ó}{{\'O}}1 {Ú}{{\'U}}1
        {ñ}{{\~n}}1 {Ñ}{{\~N}}1 {ü}{{\"u}}1 {¿}{{?`}}1 {¡}{{!`}}1,
}

% Estilo para código Python
\lstdefinestyle{python}{
    style=estilo_base,
    language=Python,
}
\renewcommand{\lstlistingname}{Código}

\newcommand{\codigo}[1]{\texttt{\small #1}}

% Comando corto para importar la solución de un ejercicio desde ejercicios/
% Uso: \codigopython{02}{Caption del ejercicio}
% Espera encontrar el archivo ejercicios/TP3_Ejercicio_#1.py
\newcommand{\codigopython}[2]{%
    \lstinputlisting[style=python, caption={#2}]{ejercicios/TP3_Ejercicio_#1.py}%
}

% ============================================================
% LISTAS
% ============================================================
\usepackage{enumitem}

% ============================================================
% ENLACES
% ============================================================
\usepackage{hyperref}
\hypersetup{
    colorlinks  = true,
    linkcolor   = black,
    urlcolor    = blue!60!black,
    citecolor   = black,
    pdftitle    = {Trabajo Práctico N° 3 - Programación II},
    pdfauthor   = {Facundo Luna},
}

% ============================================================
% ENCABEZADO Y PIE DE PÁGINA
% ============================================================
\usepackage{fancyhdr}
\pagestyle{fancy}
\fancyhf{}
\renewcommand{\headrulewidth}{0.4pt}
\renewcommand{\footrulewidth}{0.4pt}
\fancyhead[L]{\small Programación II}
\fancyhead[R]{\small Trabajo Práctico N° 3}
\fancyfoot[L]{\small Facundo Luna}
\fancyfoot[R]{\small \thepage}

% ============================================================
% DATOS DEL TRABAJO
% ============================================================

\newcommand{\universidad}{Universidad Católica de Salta}
\newcommand{\carrera}{Licenciatura en Ciencia de Datos}
\newcommand{\materia}{Programación II}
\newcommand{\trabajo}{Trabajo Práctico N° 3}
\newcommand{\profesor}{Ing. Mario Martínez}
\newcommand{\alumno}{Facundo Luna}
\newcommand{\fecha}{\today}

\newcommand{\logoincluido}{%
    \includegraphics[width=3.5cm]{../../../assets/ucasal-logo.png}\\[0.5cm]%
}
\newcommand{\logoomitido}{\vspace{1cm}}

\newcommand{\logoportada}{\logoincluido}

% ============================================================
% COMANDO VOLVER AL ÍNDICE
% ============================================================
\newcommand{\volverindice}{%
    \vspace{1em}%
    \begin{flushright}%
        \hyperlink{indice}{\small $\uparrow$ Volver al índice}%
    \end{flushright}%
}

% ============================================================
% COMANDO PARA EL ENUNCIADO DE CADA EJERCICIO
% ============================================================
\newcommand{\enunciado}[1]{%
    \begin{quote}
        \itshape #1
    \end{quote}
}

% ============================================================

\begin{document}

% ======================== PORTADA ===========================

\begin{titlepage}
    \centering

    \logoportada

    \rule{\linewidth}{1.5pt}\\[0.4cm]

    {\large\scshape \universidad}\\[0.2cm]
    {\large \carrera}\\[0.1cm]

    \rule{\linewidth}{0.5pt}\\[1.5cm]

    {\Huge\bfseries \trabajo}\\[0.8cm]
    {\LARGE \materia}\\[2cm]

    \begin{tabular}{rl}
        \textbf{Alumno:}   & \alumno   \\[0.3cm]
        \textbf{Profesor:} & \profesor \\[0.3cm]
        \textbf{Fecha:}    & \fecha    \\
    \end{tabular}

    \vfill

    \rule{\linewidth}{1.5pt}

\end{titlepage}

% ======================== ÍNDICE ============================

\newpage
\tableofcontents
\hypertarget{indice}{}
\setcounter{tocdepth}{2}

% ======================== CUERPO ============================

\newpage

\phantomsection
\section*{Aclaraciones previas}
\addcontentsline{toc}{section}{Aclaraciones previas}

El presente trabajo práctico consiste en la codificación en Python de 30 ejercicios de la Unidad correspondiente de Programación II, referidos al uso de estructuras condicionales (\texttt{if}, \texttt{elif}, \texttt{else}) y estructuras repetitivas (\texttt{for}, \texttt{while}, \texttt{range}). Cada ejercicio se resuelve en un archivo independiente, nombrado según el formato exigido por la cátedra (\texttt{TP3\_Ejercicio\_XX.py}), incluyendo como primera línea un comentario con el apellido y nombre del alumno.

El código fuente completo de cada ejercicio se encuentra disponible en el siguiente repositorio de GitHub, con el fin de facilitar su verificación y prueba:

\begin{center}
    \url{https://github.com/Facugl/data-science-coursework/tree/main/primer-anio/programacion-2/tp-3/ejercicios}
\end{center}

% ======================== EJERCICIO 1 ========================

\newpage
\section{Kilómetros por litro recorridos por un auto}

\enunciado{Escribir un programa que solicite la cantidad de kilómetros recorridos por un auto y la cantidad total de litros que consumió en ese recorrido. Luego mostrar por pantalla cuántos kilómetros hace por litro.}

\subsection*{Ejercicio en Python}

\codigopython{01}{Kilómetros por litro recorridos por un auto}

\volverindice
\newpage

% ======================== EJERCICIO 2 ========================

\section{Resta del 15\% a un número}

\enunciado{Escribir un programa que solicite un número, luego le reste el 15\%, almacenando todo en una única variable. A continuación, mostrar el resultado final por pantalla.}

\subsection*{Ejercicio en Python}

\codigopython{02}{Resta del 15\% a un número}

\volverindice
\newpage

% ======================== EJERCICIO 3 ========================

\section{Número par o impar}

\enunciado{Escribir un programa que solicite un número entero y luego muestre por pantalla si el número es par o no. (Usar función resto)}

\subsection*{Ejercicio en Python}

\codigopython{03}{Número par o impar}

\volverindice
\newpage

% ======================== EJERCICIO 4 ========================

\section{Vocal o no vocal}

\enunciado{Escribir un programa que solicite el ingreso de una letra, luego valide si es una vocal, mostrando el mensaje \enquote{Es vocal} o \enquote{No es vocal}.}

\subsection*{Ejercicio en Python}

\codigopython{04}{Vocal o no vocal}

\volverindice
\newpage

% ======================== EJERCICIO 5 ========================

\section{Suma, promedio y multiplicación de tres números}

\enunciado{Escribir un programa que solicite el ingreso de 3 números. Luego muestre por pantalla: la suma de los números, el promedio de los números y la multiplicación de los 3 números (todo en un mismo programa).}

\subsection*{Ejercicio en Python}

\codigopython{05}{Suma, promedio y multiplicación de tres números}

\volverindice
\newpage

% ======================== EJERCICIO 6 ========================

\section{Conversión de Fahrenheit a Celsius}

\enunciado{Escribir un programa que solicite al usuario el ingreso de una temperatura en escala Fahrenheit (debe permitir decimales), y luego muestre por pantalla el equivalente en grados Celsius. La fórmula de conversión que se usa para este cálculo es: Celsius = (5/9) * (Fahrenheit-32).}

\subsection*{Ejercicio en Python}

\codigopython{06}{Conversión de Fahrenheit a Celsius}

\volverindice
\newpage

% ======================== EJERCICIO 7 ========================

\section{Examen tipo múltiple choice}

\enunciado{Escribir un programa que funcione como un examen tipo múltiple choice. Pregunte al usuario: \enquote{¿Qué color se genera mezclando azul y amarillo?} y muestre las opciones \enquote{1- Rojo; 2- Violeta; 3- Verde; 4- Negro}. Luego de recibida la respuesta muestre el mensaje de \enquote{Respuesta Correcta} o \enquote{Respuesta incorrecta} según corresponda.}

\subsection*{Ejercicio en Python}

\codigopython{07}{Examen tipo múltiple choice}

\volverindice
\newpage

% ======================== EJERCICIO 8 ========================

\section{Calificación según la nota del examen}

\enunciado{Solicitar al usuario el ingreso de la nota del Examen. Luego mostrar por pantalla:}

\begin{itemize}[noitemsep]
    \item[a.] \enquote{Sobresaliente! Felicitaciones!} (nota 10)
    \item[b.] \enquote{Promocionado! Bravo!} (7 a 9)
    \item[c.] \enquote{Aprobado. Bien pero a poner más empeño!} (4 a 6)
    \item[d.] \enquote{Desaprobado. Prepararse para el recuperatorio!} (0 a 3)
\end{itemize}

\subsection*{Ejercicio en Python}

\codigopython{08}{Calificación según la nota del examen}

\volverindice
\newpage

% ======================== EJERCICIO 9 ========================

\section{Comparación y verificación de dos números}

\enunciado{Solicitar el ingreso de dos números y mostrar: cuál de los dos números es mayor; verificar si el primer número es positivo, negativo o cero; verificar si el segundo número es mayor a 0 y menor a 100.}

\subsection*{Ejercicio en Python}

\codigopython{09}{Comparación y verificación de dos números}

\volverindice
\newpage

% ======================== EJERCICIO 10 ========================

\section{Diferencia entre dos números: positiva, negativa o cero}

\enunciado{Realice un programa que solicite 2 números. Calcule la diferencia entre ellos e informe por pantalla si el resultado es positivo, negativo o cero.}

\subsection*{Ejercicio en Python}

\codigopython{10}{Diferencia entre dos números: positiva, negativa o cero}

\volverindice
\newpage

% ======================== EJERCICIO 11 ========================

\section{Paridad de tres números enteros}

\enunciado{Realice un programa que solicite el ingreso de tres números enteros, y luego en cada caso informe si el número es par o impar.}

\subsection*{Ejercicio en Python}

\codigopython{11}{Paridad de tres números enteros}

\volverindice
\newpage

% ======================== EJERCICIO 12 ========================

\section{Máxima, mínima y promedio de tres temperaturas}

\enunciado{Realice un programa que solicite ingresar tres valores de temperatura. De las temperaturas ingresadas determinar cuál de ellas es la máxima temperatura ingresada, cuál es la mínima temperatura ingresada y cuál es el promedio de las temperaturas ingresadas.}

\subsection*{Ejercicio en Python}

\codigopython{12}{Máxima, mínima y promedio de tres temperaturas}

\volverindice

% ======================== SECCIÓN BUCLES ========================

\newpage
\phantomsection
\section*{Bucles}
\addcontentsline{toc}{section}{Bucles}

A partir de este punto, los ejercicios requieren el uso de estructuras repetitivas (\texttt{for} y \texttt{while}) para resolver problemas que implican recorrer secuencias de valores, acumular resultados o repetir una acción un número determinado de veces.

% ======================== EJERCICIO 13 ========================

\newpage
\section{Números del 1 al 30 en forma ascendente}

\enunciado{Mostrar los números del 1 al 30 en forma ascendente.}

\subsection*{Ejercicio en Python}

\codigopython{13}{Números del 1 al 30 en forma ascendente}

\volverindice
\newpage

% ======================== EJERCICIO 14 ========================

\section{Números del 20 al 1 en forma descendente}

\enunciado{Mostrar los números del 20 al 1 en forma descendente.}

\subsection*{Ejercicio en Python}

\codigopython{14}{Números del 20 al 1 en forma descendente}

\volverindice
\newpage

% ======================== EJERCICIO 15 ========================

\section{Números pares entre 10 y 40}

\enunciado{Mostrar los números pares entre 10 y 40 en forma ascendente.}

\subsection*{Ejercicio en Python}

\codigopython{15}{Números pares entre 10 y 40}

\volverindice
\newpage

% ======================== EJERCICIO 16 ========================

\section{Números impares del 60 al 30 en forma descendente}

\enunciado{Mostrar los números impares del 60 al 30 en forma descendente.}

\subsection*{Ejercicio en Python}

\codigopython{16}{Números impares del 60 al 30 en forma descendente}

\volverindice
\newpage

% ======================== EJERCICIO 17 ========================

\section{Tabla de multiplicar del 6}

\enunciado{Mostrar la tabla de multiplicar del 6 (Ej. 6x1=6; 6x2=12; 6x3=18, etc. hasta 6x10).}

\subsection*{Ejercicio en Python}

\codigopython{17}{Tabla de multiplicar del 6}

\volverindice
\newpage

% ======================== EJERCICIO 18 ========================

\section{Múltiplos de 5 entre 1 y 100}

\enunciado{Mostrar los múltiplos de 5 desde el 1 al 100. (usar mod)}

\subsection*{Ejercicio en Python}

\codigopython{18}{Múltiplos de 5 entre 1 y 100}

\volverindice
\newpage

% ======================== EJERCICIO 19 ========================

\section{Números desde N hasta P}

\enunciado{Desde un número N y otro P ingresados por el usuario, mostrar todos los números desde N hasta P. (suponer que P es mayor que N)}

\subsection*{Ejercicio en Python}

\codigopython{19}{Números desde N hasta P}

\volverindice
\newpage

% ======================== EJERCICIO 20 ========================

\section{Repetición de un nombre N veces}

\enunciado{Dado el nombre de una persona y un número N ingresados por el usuario, mostrarlo N veces seguidas, tanto el nombre como su repetición. Por Ej: Ana 1, Ana 2, Ana 3.}

\subsection*{Ejercicio en Python}

\codigopython{20}{Repetición de un nombre N veces}

\volverindice
\newpage

% ======================== EJERCICIO 21 ========================

\section{Todas las tablas de multiplicar del 1 al 10}

\enunciado{Mostrar todas las tablas de multiplicar del 1 al 10, con un subtítulo entre cada tabla separando las mismas. Ej: \enquote{Tabla del 1; 1x1=1; 1x2=2; 1x3=3…1x10=10; Tabla del 2; 2x1=2; 2x2=4….}}

\subsection*{Ejercicio en Python}

\codigopython{21}{Todas las tablas de multiplicar del 1 al 10}

\volverindice
\newpage

% ======================== EJERCICIO 22 ========================

\section{Suma total y producto de 5 números}

\enunciado{Solicitar 5 números al usuario. Luego mostrar: la suma total de todos los números y el producto de todos los números.}

\subsection*{Ejercicio en Python}

\codigopython{22}{Suma total y producto de 5 números}

\volverindice
\newpage

% ======================== EJERCICIO 23 ========================

\section{Promedio de edad de un curso}

\enunciado{Dada la cantidad de alumnos de un curso ingresada por el usuario, solicitar todas las edades y mostrar el promedio general de edad del curso.}

\subsection*{Ejercicio en Python}

\codigopython{23}{Promedio de edad de un curso}

\volverindice
\newpage

% ======================== EJERCICIO 24 ========================

\section{Suma de pares y producto de impares en una lista de 10 números}

\enunciado{Dada una lista de 10 números ingresados por el usuario, calcular y mostrar la suma de todos los números pares y el producto de todos los números impares.}

\subsection*{Ejercicio en Python}

\codigopython{24}{Suma de pares y producto de impares en una lista de 10 números}

\volverindice
\newpage

% ======================== EJERCICIO 25 ========================

\section{Suma de múltiplos de 3 y producto de múltiplos de 5}

\enunciado{Dada una lista de 10 números ingresados por el usuario, calcular y mostrar la suma de todos los números múltiplos de 3 y el producto de todos los números múltiplos de 5.}

\subsection*{Ejercicio en Python}

\codigopython{25}{Suma de múltiplos de 3 y producto de múltiplos de 5}

\volverindice
\newpage

% ======================== EJERCICIO 26 ========================

\section{Cantidad de múltiplos de 4, de 2, positivos y negativos}

\enunciado{Dada una lista de 10 números ingresados por el usuario, calcular y mostrar la cantidad de números múltiplo de 4, la cantidad de números múltiplo de 2, la cantidad de números positivos y la cantidad de números negativos.}

\subsection*{Ejercicio en Python}

\codigopython{26}{Cantidad de múltiplos de 4, de 2, positivos y negativos}

\volverindice
\newpage

% ======================== EJERCICIO 27 ========================

\section{Promedio general y promedio de los números positivos}

\enunciado{Dada una lista de 10 números ingresados por el usuario, calcular y mostrar el promedio de todos los números ingresados y el promedio de los números positivos.}

\subsection*{Ejercicio en Python}

\codigopython{27}{Promedio general y promedio de los números positivos}

\volverindice
\newpage

% ======================== EJERCICIO 28 ========================

\section{Conteo de negativos y no negativos en un rango}

\enunciado{Pedir por consola dos números que representen el principio y fin de una secuencia numérica. Realizar un bucle \enquote{for} que recorra esa secuencia armada con \enquote{range} y cuente cuántos números son negativos y cuántos números son mayor o igual a cero. Tener en cuenta que \enquote{range} no incluye el número de \enquote{fin} en su secuencia, sino que va hasta el anterior.}

\subsection*{Ejercicio en Python}

\codigopython{28}{Conteo de negativos y no negativos en un rango}

\volverindice
\newpage

% ======================== EJERCICIO 29 ========================

\section{Cantidad y sumatoria de números en un rango}

\enunciado{Realice un programa que pida por consola dos números que representen el principio y fin de una secuencia numérica. Realizar un bucle \enquote{for} que recorra esa secuencia armada con \enquote{range} y cuente cuántos números ingresados hay, y la sumatoria de todos los números. Al finalizar el bucle, utilice la variable \texttt{cantidad\_numeros} y la variable \texttt{sumatoria} para calcular el promedio de todos los números ingresados. Tener en cuenta que \enquote{range} no incluye el número de \enquote{fin} en su secuencia, sino que va hasta el anterior.}

\subsection*{Ejercicio en Python}

\codigopython{29}{Cantidad y sumatoria de números en un rango}

\volverindice
\newpage

% ======================== EJERCICIO 30 ========================

\section{Calculadora con bucle y operador FIN}

\enunciado{Realice una calculadora: dentro de un bucle se debe ingresar por línea de comando dos números. Luego se ingresará como tercera entrada al programa el símbolo de la operación que se desea ejecutar: suma (+), resta (-), multiplicación (*), división (/), exponente/potencia (**).}

\enunciado{Se debe efectuar el cálculo correcto según la operación ingresada por consola. Imprimir en pantalla la operación realizada y el resultado. El programa se debe repetir dentro del bucle hasta que como operador se ingrese la palabra \enquote{FIN}, en ese momento debe terminar el programa. Se debe mostrar un mensaje de error si el operador ingresado no es alguno de los soportados o no es la palabra \enquote{FIN}.}

\subsection*{Ejercicio en Python}

\codigopython{30}{Calculadora con bucle y operador FIN}

\volverindice

\end{document}
